% !TEX program = lualatex
\documentclass[lualatex,paper=a4paper]{jlreq}
\usepackage{amsmath}
\usepackage{amssymb}
\usepackage{amsthm}
\usepackage{amsfonts}
\usepackage{mathtools}
\usepackage{graphicx}
\usepackage{multirow}
\usepackage{hyperref}
\usepackage{diffcoeff}
\usepackage{comment}
\usepackage{pdfpages}
\usepackage{float}
\usepackage[version=4]{mhchem}
\usepackage[separate-uncertainty]{siunitx}
\usepackage[math-style=ISO,warnings-off={mathtools-colon,mathtools-overbracket}]{unicode-math}
\usepackage{newunicodechar}
\usepackage{zref-clever}
\usepackage{listings}
\usepackage{xcolor}

%追加した-----------
\usepackage{tikz}
\usetikzlibrary{arrows.meta,positioning}
\usepackage{booktabs}
\usepackage{array}
%--------------------
%%%%%%%%%%%%%%%%%%%%%%%%%%%%%%%%%%%%%%%%%%%%%%%%%
\mathtoolsset{showonlyrefs,showmanualtags}
\newunicodechar{、}{，}
\newunicodechar{。}{．}
\NewDocumentCommand\、{}{\string、}
\NewDocumentCommand\。{}{\string。}
\NewDocumentCommand\degC{}{\ensuremath{^\circ\symup{C}}}
\NewDocumentCommand\abs{m}{\left|#1\right|}
\renewcommand{\lstlistingname}{コード}
%%%%%%%%%%%%%%%%%%%%%%%%%%%%%%%%%%%%%%%%%%%%%%%%%
\setmathfont{NewComputerModernMath}[
    CharacterVariant = {1}
]
\setmathfont{NewComputerModernMath}[
    range = {scr, bfscr},
    StylisticSet = 1
]
\zcDeclareLanguage[allcaps]{japanese}
\zcLanguageSetup{japanese}{
  namesep   = {\nobreak} ,
  pairsep   = {と} ,
  listsep   = {、} ,
  lastsep   = {、および} ,
  tpairsep  = {および} ,
  tlistsep  = {、} ,
  tlastsep  = {、および} ,
  notesep   = {：} ,
  rangesep  = {から} ,
  type = section ,
    Name-sg = 第 ,
    Name-pl = 第 ,
    refbounds = {,,節,} ,
  type = subsection ,
    Name-sg = 第 ,
    Name-pl = 第 ,
    refbounds = {,,小節,} ,
  type = subsubsection ,
    Name-sg = 第 ,
    Name-pl = 第 ,
    refbounds = {,,小小節,} ,
  type = paragraph ,
    Name-sg = 第 ,
    Name-pl = 第 ,
    refbounds = {,,段落,} ,
  type = appendix ,
    Name-sg = 付録 ,
    Name-pl = 付録 ,
  type = page ,
    Name-sg = p. ,
    Name-pl = pp. ,
    namesep = 〜 ,
    rangesep = \textendash ,
    rangetopair = false ,
  type = figure ,
    Name-sg = 図 ,
    Name-pl = 図 ,
  type = table ,
    Name-sg = 表 ,
    Name-pl = 表 ,
  type = item ,
    Name-sg = 項目 ,
    Name-pl = 項目 ,
  type = footnote ,
    Name-sg = 脚注 ,
    Name-pl = 脚注 ,
  type = endnote ,
    Name-sg = 後注 ,
    Name-pl = 後注 ,
  type = note ,
    Name-sg = 注記 ,
    Name-pl = 注記 ,
  type = equation ,
    Name-sg = 式 ,
    Name-pl = 式 ,
    refbounds = {,(,),} ,
  type = listing ,
    Name-sg = コード ,
    Name-pl = コード ,
}
\zcsetup{ lang = japanese, nameinlink = true }
%%%%%%%%%%%%%%%%%%%%%%%%%%%%%%%%%%%%%%%%%%%%%%%%%
\lstset{%
  basicstyle={\ttfamily},
  identifierstyle={\normalsize},
  keywordstyle={\small\bfseries\color[rgb]{0.000, 0.000, 1.000}},
  backgroundcolor={\color[gray]{0.98}},
  stringstyle={\small\ttfamily\color[rgb]{0.639, 0.082, 0.082}},
  commentstyle={\small\ttfamily\color[rgb]{0.000, 0.502, 0.000}},
  frame={tb},
  breaklines=true,
  showstringspaces=false,
  columns=fixed,
  basewidth=0.5em,
  numbers=left,
  xrightmargin=0\zw,
  xleftmargin=3\zw,
  numberstyle={\scriptsize},
  stepnumber=1,
  numbersep=1\zw,
  lineskip=-0.5ex
}
%%%%%%%%%%%%%%%%%%%%%%%%%%%%%%%%%%%%%%%%%%%%%%%%%
\jlreqsetup{
    caption_font=\mcfamily,
    caption_label_font=\mcfamily,
    caption_label_format={#1:},
    caption_after_label_space=0.5\zw
}
\ModifyHeading{section}{
    font=\normalsize\bfseries,
    lines=2,
    after_lines=0
}
\ModifyHeading{subsection}{
    font=\normalsize\bfseries,
    lines=2,
    after_lines=0
}


% \title{でんサポ}
% \author{システム提案書}
% \date{\today}


%----新しく追加-----

%\pagestyle{empty}
% 表紙の文字サイズ設定(ここの数値を変えるだけで調整できます)
\newcommand{\teamnamesize}{\fontsize{16}{22}\selectfont}
\newcommand{\teamname}{UniCo}

%sectionの文字サイズ（統一）
% 見出しの文字サイズ設定(ここを変えるだけで調整できます)
\newcommand{\sectionsize}{\fontsize{16}{24}\selectfont}
\newcommand{\subsectionsize}{\fontsize{13}{20}\selectfont}

\ModifyHeading{section}{
    font=\sectionsize\gtfamily\bfseries,
    lines=2,
    after_lines=0
}
\ModifyHeading{subsection}{
    font=\subsectionsize\gtfamily\bfseries,
    lines=2,
    after_lines=0
}

%現状の課題に対する提案色分け
\definecolor{prob}{RGB}{255,228,225}   % 課題
\definecolor{sol}{RGB}{220,235,250}    % 機能
\definecolor{eff}{RGB}{225,245,225}    % 効果
\definecolor{adm}{RGB}{240,240,240}    % 管理者

%URLの色変更
\hypersetup{
    colorlinks=true,
    linkcolor=black,     % 目次(と内部リンク全般)を黒に
    citecolor=orange,
    urlcolor=blue            % \href や \url の色
}

%itemize環境の段落ごとの調節
%\renewcommand{\labelitemi}{}                  % 1段目は記号なし
\renewcommand{\labelitemii}{}      % 2段目は記号なし
\renewcommand{\labelitemiii}{〇}    % 3段目は「–」

% 「課題1.1」のようにクリックで該当項目へ飛ぶリンクを作るマクロ
\newcommand{\issueref}[1]{\hyperref[#1]{課題\ref*{#1}}}

%------------------

\begin{document}

%---GitHubの設定---
% !TEX root = ../G-design.tex
\section{データベース設計}
\section{概要}
本レポートでは、MATLABを用いて行った全5回の音響信号処理および画像処理の実験についてまとめる。
実験では、純音やうなり、フーリエ級数などの音響信号の基礎から、画像の量子化、閾値処理、フィルタ処理、さらに2次元FFTを用いた周波数解析までを行った。
各実験で得られた波形や画像を比較し、それぞれの処理によって信号や画像がどのように変化するかを考察した。

\section{目的}
本実験の目的は、MATLABを用いて音響信号および画像データに対する基本的な信号処理を行い、それぞれの処理方法とその特徴を理解することである。
第1回では、正弦波の生成を通して周波数、振幅、位相の違いを確認し、うなり、フーリエ級数、白色ガウス雑音について学ぶ。
第2回では、FFTを用いて信号の周波数成分を解析し、周波数領域でのフィルタ処理について理解する。
第3回では、RGB画像の構造、量子化、二値化、ヒストグラム、背景差分などの基本的な画像処理を行う。
第4回では、平均化フィルタやメディアンフィルタによる雑音除去、ソーベルフィルタやラプラシアンフィルタによるエッジ検出、HSV色空間を用いた色領域の抽出を行う。
第5回では、2次元FFTを用いて画像の空間周波数を解析し、周波数領域での画像処理について理解する。
以上の実験を通して、音と画像に共通する信号処理の基本的な考え方を理解することを目的とする。
\section{方法}
\subsection{音の工学的特徴}
第1回では、\verb|Matlab|を使用して正弦波を生成し、周波数、振幅、位相を変更したときの違いを確認した。

はじめに、サンプリング周波数を16 kHz、提示時間を1秒として、440 Hzと660 Hzの純音を作成した。
作成した二つの純音をそれぞれ再生した後、波形の周期の違いを確認しやすくするため、各信号の先頭160点、すなわち約0.01秒間を取り出して同じグラフ上に表示した。
横軸を時間、縦軸を振幅とし、440 Hzと660 Hzの波形を比較した。

次に、440 Hzの純音を基準として、振幅を0.5倍および0.25倍に変更した信号を作成した。
また、初期位相を90度および180度に変更した信号を作成した。位相については、90度をπ/2、180度をπとして計算に用いた。
それぞれの条件について、元の440 Hzの信号と変更後の信号の先頭160点を同じグラフ上に表示した。
また、元の信号を左チャンネル、変更後の信号を右チャンネルとしたステレオ信号を作成し、それぞれ再生した。

うなりについては、サンプリング周波数16 kHz、提示時間4秒として、440 Hzと441 Hzの純音をそれぞれ作成した。
その後、二つの信号を加算して一つの信号とした。作成した信号を再生するとともに、数秒間の時間波形を表示し、振幅が時間とともにどのように変化するかを確認した。

フーリエ級数では、周期0.5秒の矩形波を対象とし、基本周波数を2 Hzとした。
基本周波数とその奇数倍の周波数を持つ正弦波を加算し、項数をN=1、5、25と変更した。
N=25ではfor文を用いて1次、3次、5次と奇数次の高調波を順番に加算した。
各条件で作成した波形を0秒から2秒まで同一グラフ上に表示し、項数を増加させたときの波形の変化を比較した。

最後に、サンプリング周波数16 kHz、提示時間1秒として16000点のデータを用意し、randnを用いて平均0、標準偏差0.2となる白色ガウス雑音を生成した。
生成した雑音を再生して1秒間の時間波形を表示した。
また、振幅を100区間に分けたヒストグラムを作成し、各振幅値の出現頻度を確認した。

\subsection{周波数解析}
第2回では、\verb|Matlab|のFFTを用いて時間領域の信号を周波数領域へ変換し、信号に含まれる周波数成分について確認した。

はじめに、サンプリング周波数を8192Hz、データ数を1024点として440Hzの純音を作成した。
作成した信号について、0秒から0.02秒までの時間波形を表示した。
その後、fft関数を用いてFFTを行い、得られた複素数データの絶対値を求めることで振幅スペクトルを作成した。
周波数軸はサンプリング周波数とデータ数から求め、0Hzから1000Hzまでの範囲を表示した。

次に、周波数を440.2Hzに変更し、440Hzの場合と同じくサンプリング周波数8192Hz、データ数1024点として純音を作成した。
作成した時間波形にFFTを行い、0Hzから1000Hzまでの振幅スペクトルを表示した。
440Hzと440.2Hzの結果を比較することで、信号の周波数とFFTによって得られるスペクトルとの関係を確認した。

続いて、128点の値0と128点の値1を一周期とし、それを4回繰り返すことで矩形波を作成した。
このときもサンプリング周波数は8192 Hzとした。作成した矩形波の時間波形を表示した後にFFTを行い、振幅スペクトルを求めた。
矩形波には複数の周波数成分が含まれているため、0Hzから200Hzまでの振幅スペクトルを表示し、どのような周波数にピークが現れるかを確認した。

最後に、周波数領域でのフィルタ処理を行った。
サンプリング周波数を8192Hzとし、ド、レ、ミ、ファ、ソ、ラ、シ、ドに対応する261.38Hz、293.67Hz、329.63Hz、349.23Hz、392.00Hz、440.00Hz、493.88Hz、523.23Hzの純音をそれぞれ0.25秒間生成した。
作成した八つの純音を順番に連結し、一つの時間信号とした。

連結した信号を再生して時間波形を表示した後、FFTを行って周波数領域へ変換した。
フィルタ処理前の振幅スペクトルは0Hzから600Hzまで表示した。
その後、カットオフ周波数を375Hzとするハイパスフィルタを作成した。
FFT後の正の周波数と負の周波数の両方について、絶対値が375Hz未満となる周波数成分を0とし、それ以外を1とした。
作成したフィルタをFFT後の複素数データに乗算することで、低い周波数成分を除去した。

フィルタ処理後のデータにifftを適用して時間領域へ戻し、計算結果の実部を取り出した。
最後に、フィルタ処理後の時間波形と振幅スペクトルを表示し、元の信号との違いを比較するとともに、処理後の信号を再生した。

\subsection{画像の理解}
第3回では、kut.jpgを\verb|Matlab|で読み込み、画像の画素値に対して複数の処理を行った。

最初にカラー画像から赤、緑、青の各チャネルをそれぞれ取り出し、グレイスケール画像として表示した。
その後、赤色チャネルと青色チャネルを交換したカラー画像を作成した。
これによって、RGBの各成分が元画像の色にどのように関係しているかを確認した。

次にカラー画像をグレイスケール化し、量子化ビット数を8bit、4bit、2bit、1bitに変更した。
8bit画像を基準としてビットシフトを行い、それぞれのビット数で使用できる画素値の範囲に変換した。
作成した四つの画像を並べて表示し、ビット数を減少させることで画像の見え方がどのように変化するかを確認した。

階調反転では、それぞれの量子化ビット数における最大値から各画素値を引いた。
これにより、元画像の明るい部分が暗く、暗い部分が明るくなるようにした。

閾値処理では8bitのグレイスケール画像に対して複数の閾値を設定した。
画素値が閾値より大きい場合と小さい場合に分けることで、白と黒の二値画像を作成した。

ヒストグラムでは、画素値0から255までについて、その値を持つ画素が画像内にいくつ存在するかを計算した。
この課題では\verb|hist|や\verb|histogram|を使用せず、それぞれの画素値と画像を比較し、sumを利用して個数を求めた。

最後に、人物が写っている画像と背景だけの画像を用いて背景差分を行った。
二つの画像の同じ位置にある画素同士を減算し、差が大きく現れる場所を調べた。
差分では負の値が発生する可能性があるため、uint8のまま単純に減算するのではなく、負の値を扱えるデータ型に変換してから計算した。

\subsection{画像のフィルタ処理}
第4回では、画像に雑音を加えた後に平滑化処理を行い、フィルタの種類による雑音除去効果の違いを確認した。
また、ソーベルフィルタとラプラシアンフィルタによるエッジ検出を行い、最後にHSV色空間を利用した肌色領域の検出を行った。

はじめに、kut.jpgを読み込み、RGBの各チャネルから

\[
Y = 0.3R + 0.59G + 0.11B
\]

としてグレイスケール画像を作成した。
次に、この画像にインパルス雑音と白色ガウス雑音の2種類の雑音を加えた。
インパルス雑音では、各画素に対して0から1までの一様乱数を生成し、乱数が0.01以下となった画素を黒、0.99以上となった画素を白に変更した。
これにより、画像中の約1\%を黒、約1\%を白とする雑音画像を作成した。
白色ガウス雑音では、平均0、標準偏差10の正規乱数を生成し、元のグレイスケール画像の各画素値に加算した。
作成したグレイスケール画像、インパルス雑音画像、白色ガウス雑音画像をそれぞれ表示して比較した。

次に、作成した2種類の雑音画像に対して、平均化フィルタとメディアンフィルタを適用した。
平均化フィルタでは、すべての要素が1/9である3×3のフィルタを作成し、filter2を用いて画像全体に適用した。
これにより、注目画素とその周囲8画素を含めた9画素の平均値を新しい画素値として使用した。
インパルス雑音画像と白色ガウス雑音画像の両方に同じ平均化フィルタを適用した。

メディアンフィルタでは、まず画像の周囲に1画素分のゼロパディングを行った。
その後、画像を1画素ずつ移動しながら、それぞれの画素を中心とする3×3の領域を取り出した。
取り出した9個の画素値の中央値をmedianによって求め、その値を新しい画素値として代入した。
この処理をインパルス雑音画像と白色ガウス雑音画像の両方に対して行い、平均化フィルタを適用した場合との違いを比較した。

続いて、グレイスケール画像にソーベルフィルタを適用してエッジを検出した。
使用した3×3のフィルタは、

\[
\begin{bmatrix}
-1 & 0 & 1 \\
-2  & 0  & 2  \\
-1  & 0  & 1
\end{bmatrix}
\]

および

\[
\begin{bmatrix}
-1 & -2 & -1 \\
0  & 0 & 0  \\
1  & 2 & 1
\end{bmatrix}
\]

の2種類とした。
それぞれのフィルタをfilter2によってグレイスケール画像へ適用した。
フィルタ処理によって得られた値には正負の値が含まれるため、絶対値を取り、それぞれの最大値を基準として0〜255の範囲に正規化した。
その後、水平方向ソーベルフィルタと垂直方向ソーベルフィルタの結果を画像として表示して比較した。

次に、グレイスケール画像に8近傍ラプラシアンフィルタを適用した。
使用したフィルタは、

\[
\begin{bmatrix}
1 & 1 & 1 \\
1  & -8 & 1  \\
1  & 1 & 1
\end{bmatrix}
\]

とし、filter2を用いて画像全体に適用した。
処理後の値について絶対値を求め、最大値を基準として0〜255の範囲に正規化した。

さらに、ラプラシアンフィルタ適用後の画素値についてヒストグラムを作成し、画素値の分布を確認した。
正規化した画素値の下位5\%に位置する値を下限の閾値、95\%に位置する値を上限の閾値として設定した。
以下の画素を0、上限以上の画素を255とし、その間に存在する画素については0〜255の範囲へ線形変換した。
これにより、弱い応答を抑えながら強いエッジを残した画像を作成し、閾値処理前と処理後の画像を比較した。

最後に、カラー画像face.pngを用いて肌色領域の検出を行った。
読み込んだRGB画像をrgb2hsvによってHSV色空間へ変換し、色相H、彩度S、明度Vの3成分に分けた。
肌色として判定する範囲を、Hについて0.036〜0.084、Sについて0.105〜0.592、Vについて0.577〜1.00と設定した。
3条件をすべて満たした画素を肌色領域として1、それ以外を0とするマスク画像を作成した。

作成したマスクをHSV画像のH、S、Vそれぞれに乗算し、肌色と判定された部分だけを残した。
その後、hsv2rgbを用いてRGB画像へ戻した。
元画像、肌色判定を表すマスク画像、肌色領域だけを残した画像をそれぞれ表示し、HSVの各成分に閾値を設定することで特定の色領域を抽出できるかを確認した。


\subsection{2次元フーリエ変換・逆フーリエ変換}
第5回では、\verb|Matlab|を用いて二次元画像に対するフーリエ変換を行い、画像の空間的な特徴と周波数成分との関係について確認した。
はじめに正弦波を利用した縞模様を作成して二次元FFTを行い、縞の本数や方向によってパワースペクトルがどのように変化するかを調べた。
次に、長方形を配置した画像についてフーリエ変換を行い、画像内で長方形の位置を変化させた場合のパワースペクトルを比較した。
最後に実際の画像に対して二次元FFTとローパスフィルタ処理を行い、特定の周波数成分を除去した場合に画像がどのように変化するかを確認した。

はじめに、256×256画素の画像を用意し、正弦波を利用して縞模様を作成した。
画像の列方向に沿って正弦波状に画素値を変化させることで、縦方向に伸びる縞模様を生成した。
プログラムでは縞の本数を変更できるようにし、4本、16本、64本の場合について確認できるようにした。
今回のプログラムでは64本が設定されており、縞の本数に応じて表示範囲についても変更できるようにした。

作成した縞画像にfft2を適用して二次元フーリエ変換を行い、その後fftshiftを用いて低周波成分が画像の中央に配置されるように並べ替えた。
フーリエ変換後のデータからパワースペクトルを求め、値の差を確認しやすくするために対数変換して表示した。
また、パワースペクトル全体だけでなく中央付近を拡大した画像も表示し、縞模様に対応する周波数成分がどの位置に現れるかを確認した。

次に、横方向に伸びる縞模様について調べた。
画像サイズは同じく256×256画素とし、縞の本数を4本、16本、64本の3種類とした。
先ほどとは異なり、画像の行方向に沿って正弦波状に画素値を変化させることで横縞を作成した。
3種類の画像についてそれぞれ二次元フーリエ変換を行い、低周波成分を中央へ移動させた後、パワースペクトルを対数表示した。
元画像とFFT後の結果を並べて表示することで、縞の本数を増加させた場合に周波数成分の位置がどのように変化するかを比較した。

続いて、256×256画素の黒色画像上に白色の長方形を配置した画像を作成した。
長方形は縦方向約30画素、横方向約60画素の大きさとした。
長方形の位置については、画像中央付近に配置した場合と、そこから縦方向および横方向へ60画素移動させた場合の2種類を作成した。

それぞれの長方形画像に二次元フーリエ変換を適用し、周波数成分を画像中央へ移動させた。
その後、パワースペクトルを求めるとともに、値の違いを見やすくするために対数変換したパワースペクトルも作成した。
元の長方形画像、通常のパワースペクトル、対数変換後のパワースペクトルを横に並べ、同じ形状の長方形を画像内で移動させた場合にFFTの結果がどのように変化するかを比較した。

最後に、LENNIA.bmpを読み込み、実際の画像に対する周波数領域でのフィルタ処理を行った。
読み込んだ画像に二次元フーリエ変換を適用した後、低周波成分が画像中央に配置されるように処理した。

フィルタ処理では、周波数領域の中心から一定範囲内にある成分だけを通過させる円形のフィルタを作成した。
今回のプログラムでは、中心からの距離が50以下となる領域を通過させ、それより外側の領域を除去するように設定した。
そのため、画像の低周波成分を残し、高周波成分を除去するローパスフィルタとして処理を行った。

作成したローパスフィルタをフーリエ変換後の画像に適用し、通過させる周波数成分を制限した。
その後、周波数成分の配置を元に戻して逆フーリエ変換を行い、再び画像として表示できる状態にした。
処理後のデータについて絶対値を求め、uint8型へ変換することでフィルタ処理後の画像を作成した。

最後に、元画像、フーリエ変換後の周波数成分、対数表示したパワースペクトル、作成した円形ローパスフィルタ、フィルタ適用後の周波数成分、逆フーリエ変換によって得られた画像などをそれぞれ表示した。
これらを比較することで、画像に含まれる周波数成分と、低周波成分のみを残した場合の画像の変化を確認した。


\section{結果}
\subsection{音の工学的特徴}
課題1では、440Hzと660Hzの純音を作成し、0秒から0.01秒までの波形を比較した。
図\ref{1}に示すように、どちらも振幅はおよそ−1から1の範囲で変化していたが、同じ時間内では660Hzの方が440Hzよりも多く振動していた。
440Hzでは0.01秒間に約4.4周期、660Hzでは約6.6周期の振動が含まれるため、実際のグラフでも660Hzの方が波の間隔が狭くなっていることが確認できた。

\begin{figure}[H]
    \centering
    %\includegraphics[width=0.7\linewidth]{第1回_画像/kadai1.pdf}
    \caption{440Hzと660Hzの波形}
    \label{1}
\end{figure}



同様に振幅を0.25倍にした場合には、最大振幅がおよそ0.25となった。
この場合も波形の周期や山と谷が現れる位置は元の信号と同じであり、振幅のみが小さくなっていることが確認できた。



初期位相を90度変化させた場合には、変更後の波形が元の波形に対して時間方向にずれて表示された。
グラフの開始時点を見ると、元の信号が0付近から始まっているのに対し、90度変更した信号は最大値付近から始まっている。
一方、最大振幅や波形の周期については元の信号との違いは見られなかった。

% !TEX root = ../G-design.tex
\section{機能設計１}
% !TEX root = ../G-design.tex
\section{機能設計２}
打越です
% !TEX root = ../G-design.tex
\section{はじめに}
\section{眼球運動計測実験（EyeLink）}

\subsection{概要}
% 300字程度。目的・手法・結果・主要な考察を要約する。
本実験では，2枚の顔画像から魅力的と感じる方を選択する選好判断課題を用い，
選択に至るまでの視線が意思決定にどのように関与するかを検討することを目的とした。
眼球運動計測装置EyeLink~IIを用いて課題遂行中の注視位置を記録し，
キー押しによる選択反応時刻を基準として，選択された顔を注視していた割合の時間変化を
MATLABにより解析した。その結果，選択直前に向けて選択画像への注視割合が上昇する傾向が観察された。
この結果は，視線が選好の形成過程を反映しているとする先行研究の知見と整合する。
% ↑ 300字程度になるよう、実際の結果に合わせて末尾を調整してください。

\subsection{目的}
% 300字程度以上。実験を行う理由・意義。仮説を明示する。
ヒトの意思決定は外界から得た感覚情報に基づいて形成される認知過程であり，
その過程を可視化することは行動理解において重要な課題である\cite{tohoku2023}。
とりわけ視線は，選好の形成と相互に影響し合うことが知られている。
Shimojoら\cite{shimojo2003}は，2つの顔から好ましい方を選ぶ課題において，
選択に至る直前に選択対象への注視が段階的に増加する「視線カスケード効果
（gaze cascade effect）」を報告した。この現象は，単に選好が視線を導くだけでなく，
注視そのものが選好の形成に寄与する可能性を示唆している。
そこで本実験では，選好判断課題における注視位置を計測し，
選択された画像への注視割合が選択反応時刻に近づくにつれて上昇するか否かを明らかにすることを目的とした。
仮説として，キー押し直前に向かって選択画像への注視割合が0.5を上回って増加すると予想した。

\subsection{方法}
% 400〜1000字程度以上。追試可能な水準で詳細に。
\subsubsection{実験参加者}
本実験は第2グループ（2班・11班）の学生から選ばれた20代男性１名を対象として実施した。

% 参加者の人数・年齢範囲・視力（矯正含む）などが記録されていれば追記してください。

\subsubsection{装置および刺激}
眼球運動の計測には眼球運動計測装置EyeLink~II（サンプリング周波数250~Hz）を用いた。
刺激としてディスプレイ上に2枚の顔画像を左右に並べて提示した。
被験者は各試行で魅力的と感じた方の顔をキー押しにより選択し，
その左右反応（左：52，右：54）と反応時間が記録された。
提示および反応取得の試行数は1名あたり20試行とした。

\begin{figure}[htbp]
  \centering
  \begin{subfigure}[b]{0.32\textwidth}
    \centering
    \includegraphics[width=\linewidth]{kigu.jpg}
    \caption{実験器具}
    \label{fig:eye_kigu}
  \end{subfigure}
  \hfill
  \begin{subfigure}[b]{0.32\textwidth}
    \centering
    \includegraphics[width=\linewidth]{jikkenfuukei.jpg}
    \caption{実験風景}
    \label{fig:eye_fuukei}
  \end{subfigure}
  \hfill
  \begin{subfigure}[b]{0.32\textwidth}
    \centering
    \includegraphics[width=\linewidth]{hikennsya.jpg}
    \caption{被験者の様子}
    \label{fig:eye_hikennsya}
  \end{subfigure}
  \caption{実験環境}
  \label{fig:eye}
\end{figure}
\clearpage
\subsubsection{手続き}
各試行では，画面中央に注視点を提示した後に2枚の顔画像を左右に提示し，
被験者は魅力的と判断した側の顔に対応するキーを押した。
選択反応（キー押し）の時刻を0~msとし，その1秒前（$-996$~msから0~msまで，4~ms間隔）の
区間を解析対象とした。

\subsubsection{データ解析}
解析にはMATLABを用いた。各試行のEyeLinkデータ（時間，測定眼情報，画面状態，
左目注視位置$x$座標等を含むCSVファイル）を\texttt{readmatrix}関数で読み込んだ。
計測できなかった時刻は欠損値として扱い，\texttt{ismember}関数を用いて
サンプリング時刻ごとにデータの有無を判定した。注視点提示時（画面状態1）の
$x$座標の平均を画面中心とみなし，画像提示時の$x$座標がこの中心より小さい場合を左（52），
大きい場合を右（54）と判定した。この注視方向がキー出力データの選択方向と一致した場合を1，
一致しなかった場合を0，欠損を$-1$として各試行の行列（\texttt{expos}）にまとめた。
続いて20試行分を統合し，各サンプリング時刻について欠損値を除外したうえで，
選択画像を注視していた割合（反応率）を算出した。
横軸を刺激提示（キー押し）からの時間，縦軸を選択画像に対する注視確率とするグラフを作成した。

% 追加課題（全班データの統合平均・標準誤差、検定）を行った場合はここに追記してください。

\subsection{結果}
% 結果と考察で800〜1000字程度以上。ここでは事実のみを記述する。
選択画像に対する注視確率の時間変化を図\ref{fig:eye_graph}に示す。

\begin{figure}[htbp]
  \centering
  \includegraphics[width=0.8\hsize]{graph.pdf}
  \caption{注視確率の時間変化}
  \label{fig:eye_graph}
\end{figure}
キー押しの約1秒前の時点では注視確率は約0.5であった。その後、キー押し前およそ800〜550msの区間では一時的に0.35前後まで低下したが、キー押し前約500ms以降は選択反応時刻に近づくにつれて単調に上昇し、キー押し直前には約0.75に達した。
\subsection{考察}
% 結果に基づき、根拠を示しながら推論する。事実と推測を分ける。
本実験では，選択反応時刻に近づくにつれて選択画像への注視確率が上昇する傾向が得られた。
これはShimojoら\cite{shimojo2003}が報告した視線カスケード効果と一致する結果であると考えられる。
選択直前に注視が特定の対象へ収束する現象は，視線が単に選好の結果を反映するだけでなく，
注視の継続自体が対象への選好を高める方向に働いている可能性を示唆する。
また，松宮ら\cite{tohoku2023}は意思決定の信号が眼球運動系に連続的に流れることを示しており，
本結果における注視確率の漸増も，意思決定過程が眼球運動に持続的に反映された結果として解釈できる。
% ↓ 推測であることを明示。
一方、キー押し前800〜550msの区間で注視確率が一時的に0.5を下回った点については、選択が確定する以前の段階では両画像を交互に比較する探索的な視線移動が生じ、選択される側とは逆の画像へ視線が向かう時間が相対的に多かった可能性が考えられる。その後の注視確率の上昇は、比較を経て選好が形成されるにつれ視線が選択対象へ収束していく過程を反映していると推測される。
なお，キー押し直後ではなく直前から注視確率が高まった点については，
選好が形成される前段階から視線の偏りが生じていた可能性が推測される。
本実験の限界として，試行数が20試行と限られ，被験者数も少ないため，
個人差や刺激画像の偏りの影響を排除しきれていない点が挙げられる。
今後は被験者数と試行数を増やし，刺激画像の魅力度を統制したうえで検証することが課題である。
% !TEX root = ../G-design.tex
\section{保守・運用}
% !TEX root = ../G-design.tex
\section{ネットワーク設計}
% !TEX root = ../G-design.tex
\section{セキュリティ対策}
セキュリティ対策の項目として以下が考えられる.
\begin{itemize}
    \item ファイアウォールや侵入検知システム（IDS/IPS）による外部からの不正アクセス防止を行う.

    \item Webサイトの常時SSL化（HTTPS対応）により通信を暗号化する.

    \item 氏名やプロフィール等の個人情報を適切に管理し, パスワードはハッシュ化して保存する.

    \item OS・ミドルウェア・アプリケーションの脆弱性対策を定期的に実施する.

    \item 管理者アカウントに対して強固なパスワードポリシーや多要素認証を導入する.

    \item 投稿内容に個人情報や外部の不審なリンクが含まれていないかを確認し, フィッシング等の被害防止に努める.

    \item 利用者が詐欺被害等に遭わないよう, 通報機能や注意喚起を通じて安全な利用環境の維持を図る.
\end{itemize}
% !TEX root = ../G-design.tex
\section{UI設計}
aaaaaaa
\section{結果}
UDPトラフィックを用いて、帯域幅を0.5Mbから10Mbまで0.5Mbずつ増やしていき、
それぞれのスループットとドロップ率を測定した。

以下にその結果の表およびグラフを以下に示す。

\begin{table}[H]
\centering
\caption{帯域幅ごとのスループットと損失率}
\begin{tabular}{|c|c|c|}
\hline
Bandwidth (Mb) & Throughput (Mb) & Drop Rate \\
\hline
0.5 & 0.4992 & 0.94 \\
1.0 & 0.998 & 0.87 \\
1.5 & 1.498 & 0.81 \\
2.0 & 1.9972 & 0.75 \\
2.5 & 2.4964 & 0.69 \\
3.0 & 2.996 & 0.62 \\
3.5 & 3.4952 & 0.56 \\
4.0 & 3.9944 & 0.50 \\
4.5 & 4.494 & 0.44 \\
5.0 & 4.9932 & 0.37 \\
5.5 & 5.4924 & 0.31 \\
6.0 & 5.992 & 0.25 \\
6.5 & 6.4912 & 0.19 \\
7.0 & 6.9908 & 0.12 \\
7.5 & 7.49 & 0.06 \\
8.0 & 7.9892 & 0.00 \\
8.5 & 7.9892 & 0.00 \\
9.0 & 7.9892 & 0.00 \\
9.5 & 7.9892 & 0.00 \\
10.0 & 7.9892 & 0.00 \\
\hline
\end{tabular}
\end{table}

\subsection{各帯域幅とスループット}

\begin{figure}[H]
    \centering
    \includegraphics[width=\textwidth]{Throughput.pdf}
    \caption{帯域幅とスループット}
\end{figure}
% !TEX root = ../G-design.tex
\section{利用の流れ}
%-----------------

\begin{titlepage}
  \centering
  \vspace*{35mm}
  % 日本語タイトル
  {\fontsize{30}{38}\selectfont\gtfamily\bfseries システム提案書\par}
  \vspace{2mm}
  % 英語名(欧文ゴシック・太字)
  %{\fontsize{30}{36}\selectfont\sffamily\bfseries SceneTrip\par}
  \vspace{30mm}
  % 文書名
  {\fontsize{50}{34}\selectfont\gtfamily\bfseries でんサポ\par}
  \vspace{45mm}
  % 版数
  {\fontsize{24}{24}\selectfont\gtfamily  第1.0版\par}
  \vspace{10mm}
  % 日付
  {\fontsize{20}{20}\selectfont\gtfamily\today\par}
  \vspace{15mm}
  % チーム名(ゴシック体)
  {\teamnamesize\gtfamily\teamname\par}
  \vfill
\end{titlepage}

\setcounter{page}{1}
\pagestyle{plain}
% ここから 1 ページ目として数える

% --- 全体の総合タイトル
%\maketitle

% --- 目次の自動生成 ---
{\fontsize{14}{24}\selectfont\gtfamily
\tableofcontents
}
\clearpage



\phantomsection
\section*{参考文献}
\addcontentsline{toc}{section}{参考文献}

\begin{enumerate}
  %著者名.ホームページ名.URL.閲覧日.
  %高齢社会白書（現状の課題文中：冒頭）
  \item \hypertarget{ref:kourei2026}{内閣府本府．\href{https://www8.cao.go.jp/kourei/whitepaper/w-2026/html/zenbun/s1_1_1.html}{令和8年版高齢社会白書（全体版）}．閲覧日：令和8年10/15}
  %情報通信白書|インターネット（現状の課題：図１,３）
  \item \hypertarget{ref:soumu2026}{総務省，\href{https://www.soumu.go.jp/johotsusintokei/whitepaper/ja/r08/html/nd21b120.html}{令和8年度版 情報通信白書（年次報告書）｜インターネット}．閲覧日：令和8年10/15}
  %
  \item \hypertarget{ref:soumu2021}{総務省，\href{https://www.soumu.go.jp/johotsusintokei/whitepaper/ja/r03/html/nd111430.html}{令和3年度版 情報通信白書（年次報告書）｜デジタル活動支援}．閲覧日：令和8年10/15}
  %情報通信白書|デジタル活用の動向（現状の課題：図２）
  \item \hypertarget{ref:soumu2026b}{総務省，\href{https://www.soumu.go.jp/johotsusintokei/whitepaper/ja/r08/html/nd21b110.html}{令和8年度版 情報通信白書（年次報告書）｜情報通信機器・端末}．閲覧日：令和8年10/15}
\end{enumerate}

\end{document}